\title{\pkg{amore}: A Unified Case-Control Partial-Likelihood Toolchain for
Simulating, Fitting, and Checking Relational and Hyper-Event Models in \proglang{R}}

\author{
  Francisco Richter \\ Universit\`a della Svizzera italiana
  \And
  Martina Boschi \\ Universit\`a della Svizzera italiana
  \AND
  Melania Lembo \\ Universit\`a della Svizzera italiana
  \And
  Ernst C. Wit \\ Universit\`a della Svizzera italiana
}

\Plainauthor{Francisco Richter, Martina Boschi, Melania Lembo, Ernst C. Wit}
\Plaintitle{amore: A Unified Case-Control Partial-Likelihood Toolchain for
Simulating, Fitting, and Checking Relational and Hyper-Event Models in R}

\date{}

\begin{abstract}
Relational event models describe how interactions in a network unfold in
continuous time, with intensities driven by the history of past events.
Recent methodological advances --- timing and closure variants of the
endogenous statistics, smooth and time-varying effects, global covariates,
hyper-events, and martingale-residual goodness-of-fit tests --- have so far
lived in separate implementations, making complete analyses difficult to
carry out and to reproduce. We present \pkg{amore}, an \proglang{R} package
that brings the full workflow --- simulation, case-control sampling,
endogenous-statistic construction, estimation, and model checking --- under
a single case-control partial likelihood. The estimation framework offers a
spectrum of intensity specifications behind one interface: a linear
conditional logit, smooth additive effects --- fitted in memory or, at scale,
by stochastic gradient over per-covariate spline bases --- and a neural
network scoring the full covariate vector, trading interpretability against
flexibility in a controlled way. A generative engine simulates from the
same model class it estimates, which yields an internal consistency check:
statistics computed during simulation and reconstructed from the event
history agree to numerical precision, and effects planted in simulation are
recovered in estimation. We illustrate the workflow on classroom,
e-mail, and social-network event data.
\end{abstract}

\Keywords{relational event models, partial likelihood, case-control sampling,
network event data, goodness of fit, \proglang{R}}
\Plainkeywords{relational event models, partial likelihood, case-control
sampling, network event data, goodness of fit, R}
